\documentclass[line]{resonancecv}

\cvtitle{Enes Kemal Ergin -- Resonance CV}
\name{Enes Kemal Ergin, PhD}
\tagline{Computational biologist and bioinformatician}

\contactline{\textit{Cell:} 778-751-1649}{\textit{GitHub:} \href{https://github.com/eneskemalergin}{github.com/eneskemalergin}}
\contactline{\textit{Work:} 604-875-2000 ext:7146}{\textit{LinkedIn:} \href{https://www.linkedin.com/in/eneskemalergin/}{linkedin.com/in/eneskemalergin}}
\contactline{\textit{E-mail:} \href{mailto:eergin@bcchr.ca}{eergin@bcchr.ca}}{\textit{Portfolio:} \href{https://eneskemalergin.github.io}{eneskemalergin.github.io}}
\contactline{\href{mailto:eneskemalergin@gmail.com}{eneskemalergin@gmail.com}}{\textit{Google Scholar:} \href{https://scholar.google.com/citations?user=6iZzHecAAAAJ\&hl=en}{Enes K. Ergin}}

\begin{document}

\begin{resume}

\section{Profile}
Computational biologist and bioinformatician with expertise in high-throughput omics, cancer proteomics, statistical method development, and machine learning for biomedical research.

\section{Education}
\degreeentry{University of British Columbia (UBC), Vancouver, BC Canada}{PhD, Bioinformatics}{May 2019 -- Apr 2025}{%
\begin{list2}
\item \textbf{Advisor:} Philipp F. Lange
\item \textbf{Thesis:} \href{https://open.library.ubc.ca/cIRcle/collections/ubctheses/24/items/1.0448334}{``Computational interrogation of proteoform dynamics in pediatric cancer''}
\end{list2}}

\degreeentry{University of British Columbia (UBC), Vancouver, BC Canada}{MSc, Bioinformatics}{Sep 2017 -- May 2019}{%
Transferred directly into the PhD program.}

\degreeentry{North American University (NAU), Houston, Texas USA}{BSc, Computer Science (Software Engineering)}{Sep 2013 -- May 2017}{%
\begin{list2}
\item \textbf{GPA:} 3.72/4.00
\item \textbf{Honors:} Exceptional Merit Scholarship Recipient, 2012 -- 2017
\end{list2}}

\section{Experience}
\cventry{Staff Bioinformatician}{BC Children's Hospital Research Institute}{Jun 2025 -- Present}{%
\begin{list2}
\item Develop and maintain computational infrastructure for large-scale proteomics data analysis across multiple childhood cancer cohorts.
\item Implement statistical methods and machine learning approaches to support translational research in pediatric oncology.
\end{list2}}

\cventry{Graduate Research Assistant}{BC Children's Hospital Research Institute}{Sep 2017 -- Apr 2025}{%
\begin{list2}
\item Developed and implemented \href{https://github.com/LangeLab/Analysis_of_QuEStVar_Manuscript}{QuEStVar}, expanding statistical explainability through equivalence testing and leading to a first-author publication in \emph{Journal of Proteome Research}.
\item Built \href{https://github.com/LangeLab/SQuAPP}{SQuAPP}, a tool to improve peptide-level identification and quantification in proteomics data, published in \emph{Bioinformatics}.
\item Created \href{https://github.com/LangeLab/ProteoForge_Analysis}{ProteoForge}, an imputation-aware framework for differential proteoform discovery in Neuroblastoma.
\end{list2}}

\cventry{Visiting Researcher}{Harvard Medical School}{Nov 2015 -- Sep 2016}{%
\begin{list2}
\item Predicted determinants of alternative mRNA splicing using a deep convolutional neural network trained on 11 histone modification datasets.
\item Developed a multi-task learning CNN for transcription factor binding-site prediction across cell types and participated in the ENCODE-DREAM challenge.
\item \textbf{Tools:} Lua/Torch, TensorFlow, Python (Pandas, NumPy, Matplotlib), Linux shell tools.
\end{list2}}

\cventry{Undergraduate Research Assistant}{North American University}{Sep 2015 -- May 2017}{%
\begin{list2}
\item Led student activity in the Bioinformatics Lab and mentored four junior and sophomore students.
\item Tutored newcomers with prepared tasks and facilitated weekly presentations on bioinformatics topics.
\item Designed an open source bioinformatics curriculum with Open Source Society on GitHub.
\end{list2}}

\cventry{Undergraduate Research Assistant}{Yeditepe University}{Aug 2014 -- May 2015}{%
\begin{list2}
\item Worked remotely as a genomic data scientist investigating candidate inhibitor molecules.
\item Queried NCBI PubMed and used virtual docking tools to evaluate molecular targets and likely inhibitors.
\end{list2}}

\cventry{Undergraduate Research Assistant}{Texas Institute of Education and Research (TIBER)}{Sep 2014 -- Apr 2015}{%
\begin{list2}
\item Worked as an experimental biologist on sinusitis bacteria drug development and testing.
\item Prepared agar solutions and bacterial cultures, conducted liquid and solid drug tests, and maintained experimental records.
\end{list2}}

\section{Teaching}
\cventry{Teaching Assistant}{North American University}{Sep 2015 -- May 2017}{%
Co-taught four undergraduate computer science courses for classes of more than 60 students and shared responsibility for labs, supplementary sessions, exams, homework, and grading.
\begin{list2}
\item COMP 3317 Algorithms
\item COMP 3320 Programming Languages
\item COMP 3322 Software Engineering
\item COMP 2415 Systems Programming
\end{list2}}

\cventry{Mathematics and Computer Science Tutor}{North American University}{Jan 2017 -- May 2017}{%
Tutored students in subjects ranging from college algebra and differential equations to CS1 and Data Mining.}

\cventry{Instructor}{North American University}{Feb 2015 -- Apr 2015}{%
Taught introductory Python programming to 35 students, faculty, and staff; this was the first course taught solely by a student at NAU. (\href{https://github.com/NAU-Python-Class/Py101-Spring-15}{GitHub})}

\section{Publications}
\cvsubheading{Manuscripts Under Review}
\pubentry{\textbf{Ergin, E. K.}, Conrrero, A., Ferguson, K. M., Lange, P. F. (2025). ProteoForge: An Imputation-Aware Framework for Differential Proteoform Discovery in Bottom-Up Proteomics, \emph{bioRxiv}, under review at \emph{Journal of Proteome Research}.}

\cvsubheading{Peer-Reviewed Articles and Thesis}
\pubentry{\textbf{Ergin, E. K.} (2025). \href{https://open.library.ubc.ca/cIRcle/collections/ubctheses/24/items/1.0448334}{Computational interrogation of proteoform dynamics in pediatric cancer}. PhD thesis, University of British Columbia, Vancouver, BC, Canada. Supervisor: Philipp F. Lange.}

\pubentry{Barnabas, G. D., Bhat, T. A., Goebeler, V., Leclair, P., Azzam, N., Melong, N., Anderson, C., Gom, A., An, S., \textbf{Ergin, E. K.}, Shen, Y., Conrrero, A., Mungall, A. J., Mungall, K. L., Maxwell, C. A., Reid, G. S. D., Hirst, M., Jones, S., Chan, J. A., Senger, D. L., Berman, J. N., Parker, S. J., Bush, J. W., Strahlendorf, C., Deyell, R. J., Lim, C. J., Lange, P. F. (2025). Proteomics and personalized PDX models identify treatment for a progressive malignancy within an actionable timeframe, \href{https://pmc.ncbi.nlm.nih.gov/articles/PMC11982353/}{doi: 10.1038/s44321-025-00107-w}, \emph{EMBO Molecular Medicine}.}

\pubentry{\textbf{Ergin, E. K.}, Myung, J. J. K., Lange, P. F. (2024). Statistical Testing for Protein Equivalence Identifies Core Functional Modules Conserved across 360 Cancer Cell Lines and Presents a General Approach to Investigating Biological Systems, \href{https://doi.org/10.1021/acs.jproteome.4c00131}{doi: 10.1021/acs.jproteome.4c00131}, \emph{Journal of Proteome Research}.}

\pubentry{Lorentzian, A. C., Rever, J., \textbf{Ergin, E. K.}, Guo, M., Akella, N. M., Rolf, N., Lim, C. J., Reid, G. S. D., Maxwell, C. A., Lange, P. F. (2023). Targetable lesions and proteomes predict therapy sensitivity through disease evolution in pediatric acute lymphoblastic leukemia, \href{https://doi.org/10.1038/s41467-023-42701-9}{doi: 10.1038/s41467-023-42701-9}, \emph{Nature Communications}.}

\pubentry{\textbf{Ergin, E. K.}, Uzozie, A. C., Chen, S., Su, Y., Lange, P. F. (2022). SQuAPP: Simple Quantitative Analysis of Proteins and PTMs, \href{https://doi.org/10.1093/bioinformatics/btac628}{doi: 10.1093/bioinformatics/btac628}, \emph{Bioinformatics}.}

\pubentry{Uzozie, A. C., \textbf{Ergin, E. K.}, Rolf, N., Tsui, J., Lorentzian, A., Weng, S. H., Nierves, L., Smith, T., Lim, C. J., Maxwell, C., Reid, G. S. D., Lange, P. F. (2022). PDX models reflect the proteome landscape of pediatric acute lymphoblastic leukemia but divert in select pathways, \href{https://dx.doi.org/10.1186\%2Fs13046-021-01835-8}{doi: 10.1186/s13046-021-01835-8}, \emph{Journal of Experimental and Clinical Cancer Research}.}

\pubentry{Weng, S. H.*, Demir, F.*, \textbf{Ergin, E. K.}, Dirnberger, S., Uzozie, A., Tuscher, D., Nierves, L., Tsui, J., Huesgen, P. F., Lange, P. F. (2019). Sensitive determination of proteolytic proteoforms in limited microscale proteome samples, \href{https://doi.org/10.1074/mcp.TIR119.001560}{doi: 10.1074/mcp.TIR119.001560}, \emph{Molecular and Cellular Proteomics}.}

\pubentry{Kocabas, F., \textbf{Ergin, E. K.} (2016). Identification of small molecule binding pocket for inhibition of Crimean-Congo hemorrhagic fever virus OTU protease, \href{https://doi.org/10.3906/biy-1501-56}{doi: 10.3906/biy-1501-56}, \emph{Turkish Journal of Biology}.}

\section{Software}
\cvdateditem{ProteoForge}{2025}{Imputation-aware framework for differential proteoform discovery from bottom-up proteomics data.\\
\href{https://github.com/LangeLab/ProteoForge_Analysis}{GitHub} \textbar{} \href{https://www.biorxiv.org/content/10.64898/2025.12.12.694008}{Preprint}}

\cvdateditem{QuEStVar}{2024}{R package for statistical equivalence testing in proteomics and conserved functional module discovery.\\
\href{https://github.com/LangeLab/Analysis_of_QuEStVar_Manuscript}{GitHub} \textbar{} \href{https://doi.org/10.1021/acs.jproteome.4c00131}{Publication}}

\cvdateditem{SQuAPP}{2022}{R Shiny application for end-to-end quantitative proteomics and PTM analysis.\\
\href{https://github.com/LangeLab/SQuAPP}{GitHub} \textbar{} \href{https://doi.org/10.1093/bioinformatics/btac628}{Publication}}

\section{Presentations}
\cvsubheading{Oral Presentations}
\pubentry{\textbf{Ergin, E. K.} (2023). Exploring Stable Proteome in Cancer Cell Line Atlas with QuEStVar, Trainee Omics Group (TOG), Jun 2023.}

\pubentry{\textbf{Ergin, E. K.} (2023). InPACCT: Integrated Proteomics Analysis of Curated Childhood Tumours, Canadian National Proteomics Network, May 2023.}

\pubentry{\textbf{Ergin, E. K.} (2020). Normalization in Proteomics: Past, Present, and Future, Trainee Omics Group (TOG), May 2020.}

\cvsubheading{Workshop Presentations}
\pubentry{\textbf{Ergin, E. K.} (2023). Proteomics Workshop -- Common Downstream Analysis Pipeline, Trainee Omics Group (TOG), Jul 2023.}

\cvsubheading{Poster Presentations}
\pubentry{Jenane, L., Nierves, L., \textbf{Ergin, E. K.}, Lange, P. F. (2024). Characterization of a New Type of Neoantigen in T-Cell Acute Lymphoblastic Leukemia (T-ALL) by Cell Surface Terminomics, US HUPO 2024, Portland, OR, USA.}

\pubentry{Yadegari, M., \textbf{Ergin, E. K.}, Lange, P. F. (2023). Detection of Shed Surface Proteins and Processed Oncoproteoforms in Biofluids, Canadian Cancer Research Conference 2023, Halifax, NS, Canada.}

\pubentry{\textbf{Ergin, E. K.}, Rostin, K., Lange, P. F. (2023). InPACCT: Integrated proteomics analyses of curated childhood tumours, Canadian National Proteomics Network, Regina, SK, Canada.}

\pubentry{\textbf{Ergin, E. K.}, Lange, P. F. (2023). Stable Proteome in Cancer Cell Lines with Combined Testing, Canadian National Proteomics Network, Regina, SK, Canada.}

\pubentry{Rostin, K., \textbf{Ergin, E. K.}, Lange, P. F. (2023). Insightful, Integrated Analyses of Public Pediatric Cancer Proteomics with InPACCT, BIG 23 Research Day, Vancouver, BC, Canada.}

\pubentry{Weng, S. H., Demir, F., \textbf{Ergin, E. K.}, Dirnberger, S., Uzozie, A., Tuscher, D., Nierves, L., Tsui, J., Huesgen, P. F., Lange, P. F. (2019). Sensitive determination of proteolytic proteoforms in limited samples, 2019 Pathology Day, Vancouver, BC, Canada.}

\pubentry{Ji, J., Thibodeau, M., \textbf{Ergin, E. K.}, Culibrk, L., Smith, T. (2018). SWI/SNF Chromatin Remodeling Complex in Clear Cell Ovarian Cancer. Stat540 Group Project Poster Session.}

\section{Projects}
\cvdateditem{Open Source Bioinformatics Curriculum}{Jun 2016}{Developed an open source bioinformatics curriculum with contributions from the open source community.\\
\href{https://github.com/open-source-society/bioinformatics}{GitHub}}

\cvdateditem{Scholar Development Center}{Feb 2016}{Created a non-profit community organization under the Raindrop Foundation to support Turkish undergraduate students in Texas pursuing academia or industry.}

\cvdateditem{Essential Algorithms in Python}{Sep 2014 -- Present}{Implemented algorithms from \emph{A Practical Approach to Computer Algorithms} by Rod Stephens and extended the repository with additional algorithms and number-theory implementations.}

\section{Awards}
\cvminorrow{Four Year Doctoral Fellowship Award}{2019 -- 2023}
\cvminorrow{Michael Cuccione Childhood Cancer Foundation Graduate Studentship}{2018 -- 2019}
\cvminorrow{North American University Exceptional Merit Scholarship}{2012 -- 2017}
\cvminorrow{North American University President's Honor Roll}{2015 -- 2017}
\cvminorrow{North American University Outstanding Student of the Year}{2015}

\section{Memberships}
\cvminorrow{ISCB (International Society of Computational Biology)}{Aug 2016 -- Present}
\cvminorrow{ACM (Association for Computing Machinery)}{Feb 2013 -- Present}

\section{Leadership}
\cvminorrow{NAU Future Leaders Club, Founder and Director}{Apr 2015 -- May 2017}
\cvminorrow{Student Government, VP of Unity and Social Justice}{Sep 2014 -- May 2017}
\cvminorrow{Student Government, VP of Finance}{Nov 2014 -- Apr 2015}
\cvminorrow{NAU ACM, Vice President}{Sep 2015 -- May 2016}
\cvminorrow{NAU ACM, Secretary}{Feb 2013 -- Sep 2014}
\cvminorrow{NAU Kazakh Student Association, Club Advisor}{Sep 2016 -- May 2017}

\section{Skills}
\textbullet\ \textbf{Languages and Libraries:} Python (NumPy, Pandas, Polars, Matplotlib, Seaborn, SciPy, Biopython, scikit-learn, PyTorch, TensorFlow), R (tidyverse, shiny, limma, DESeq2, MSstats), Java, LaTeX, C/C++, Bash, JavaScript, HTML, CSS\par
\textbullet\ \textbf{Data and Systems:} SQL, MySQL, MongoDB, Unix/Linux, macOS, Windows

\section{Certificates}
\cvminorrow{Neural Networks and Deep Learning, Coursera}{Nov 2017}
\cvminorrow{Statistics for Genomic Data Science, With Distinction, Coursera}{Nov 2016}
\cvminorrow{Command Line Tools for Genomics Data Science, With Distinction, Coursera}{Oct 2016}
\cvminorrow{Machine Learning, Coursera}{Aug 2016}
\cvminorrow{Bioconductor for Genomic Data Science, With Distinction, Coursera}{Jul 2016}
\cvminorrow{Algorithms for DNA Sequencing, With Distinction, Coursera}{May 2016}
\cvminorrow{Bioinformatics Methods I and II, With Distinction, Coursera}{Dec 2015}
\cvminorrow{Introduction to Genomic Technologies, With Distinction, Coursera}{Oct 2015}
\cvminorrow{Python for Genomic Data Science, With Distinction, Coursera}{Oct 2015}
\cvminorrow{Introduction to R, Honor Certificate, EdX}{Sep 2015}
\cvminorrow{Getting and Cleaning Data, Coursera}{Mar 2015}
\cvminorrow{R Programming, With Distinction, Coursera}{Feb 2015}
\cvminorrow{Introduction to R, DataCamp}{Jan 2015}
\cvminorrow{Data Scientist's Toolbox, With Distinction, Coursera}{Dec 2014}
\cvminorrow{Programming for Everybody (Python), With Distinction, Coursera}{Jun 2014}

\end{resume}

\end{document}